% 00_formatting.tex — требования к оформлению всех документов команды.
% Документ сам оформлен по этим требованиям. Собрать: ./build.sh docs/00_formatting.tex
\documentclass{tpks}
\doctitle{Требования к оформлению документов}
\docversion{2}
\docdate{20.11.2026}

\begin{document}
\makeheader
\begin{changelog}
\change{1}{20.11.2026}{Первая версия: титульный лист, лист регистрации изменений в начале документа, выделение изменений относительно принятой версии, шаблоны документов}{Чернов А. В.}{Замечания к структуре документов}
\change{2}{20.11.2026}{Запрещены ссылки на слайды, конспект и рабочие материалы; сводная таблица идентификаторов в документах с требованиями; служебные символы «; \#версия» в заголовке и ссылке, порядок объявления требований}{Чернов А. В.}{Замечания преподавателя}
\end{changelog}

\section{Введение}

\subsection{Назначение}

Документ задаёт единое оформление всех документов команды: пяти документов седьмого семестра и документов, которые появятся в восьмом. Если правило здесь расходится с образцом прошлых лет, действует это правило.

\subsection{Область применения}

Оформление задаёт шаблон \texttt{docs/tpks.cls}. Большую часть правил он выполняет сам: строит титульный лист, нумерует разделы, таблицы и рисунки, печатает колонтитулы. Остальное зависит от автора. В таблице~\ref{tab:check} раздела~\ref{sec:check} отмечено, какие правила проверяет автор перед сдачей.

\subsection{Термины и сокращения}

\textbf{Принятая версия} — версия документа, которую сдали преподавателю и которая принята. Относительно неё выделяются изменения.

\textbf{Экземпляр с выделенными изменениями} — PDF документа, где показано всё, что изменилось после принятой версии.

\textbf{Шаблон документа} — заготовка документа одного типа: требования, план или стандарт. Новый документ начинают с копии шаблона.

\subsection{Ссылки}

\begin{enumerate}
  \item Замечания преподавателя к структуре документов: титульный лист, лист регистрации в начале документа, выделение изменений относительно принятой версии, единые шаблоны, отказ от ссылок на слайды.
\end{enumerate}

\section{Состав и порядок частей}

Каждый документ состоит из частей в порядке таблицы~\ref{tab:parts}. Порядок проверяет шаблон: без титульного листа или листа регистрации, а также с листом регистрации после первого раздела документ не соберётся.

\begin{tabl}{|c|P{4.2cm}|L|P{3.2cm}|}{Части документа}[tab:parts]{\textbf{№} & \textbf{Часть} & \textbf{Содержание} & \textbf{Как получить}}
1 & Титульный лист & Университет и курс, название документа, проект, версия и дата, состав команды, преподаватель & \texttt{\textbackslash makeheader} \\ \hline
2 & Лист регистрации изменений & Все версии документа, начиная с первой & окружение \texttt{changelog} \\ \hline
3 & Содержание & Разделы трёх уровней с номерами страниц & печатается само после листа регистрации \\ \hline
4 & Основной текст & Разделы по шаблону документа (раздел~\ref{sec:templates}) & пишет автор \\ \hline
5 & Приложения & Иллюстрации и справочные материалы, которые не являются требованиями & последние разделы, заголовок начинается со слова «Приложение» \\ \hline
\end{tabl}

\section{Титульный лист}

Титульный лист — первая страница. Номер страницы на нём не печатается, но учитывается: лист регистрации начинается со страницы 2.

Название документа, версию и дату автор задаёт командами \texttt{\textbackslash doctitle}, \texttt{\textbackslash docversion} и \texttt{\textbackslash docdate}. Проект, курс, преподаватель и состав команды берутся из \texttt{docs/project.tex}, в документе их не повторяют. Если документ писала часть команды, авторов указывают командой \texttt{\textbackslash docauthors}: строка «Авторы документа» появится под версией.

\section{Лист регистрации изменений}

Лист идёт сразу после титульного листа и до содержания. В нём пять столбцов (таблица~\ref{tab:changelog}).

\begin{tabl}{|P{3.2cm}|L|}{Столбцы листа регистрации изменений}[tab:changelog]{\textbf{Столбец} & \textbf{Что писать}}
Версия & Целое число. Первая версия — 1, каждая следующая на 1 больше \\ \hline
Дата & Дата версии в виде ДД.ММ.ГГГГ \\ \hline
Что изменено & Коротко по существу: какие разделы и какие требования (по ID) изменены \\ \hline
Кто & Фамилия и инициалы автора изменения \\ \hline
Основание & «Задание к паре ДД.ММ.ГГГГ», «Замечание преподавателя от ДД.ММ.ГГГГ» или номер ЗИ \\ \hline
\end{tabl}

Правила:
\begin{itemize}
  \item одна строка на версию, строки идут по возрастанию версий;
  \item строки прежних версий не меняют и не удаляют;
  \item \texttt{\textbackslash docversion} и \texttt{\textbackslash docdate} совпадают с последней строкой листа;
  \item сведения о версиях и датах правок пишут только в листе регистрации, а не в тексте разделов.
\end{itemize}

\section{Страница и текст}

Параметры страницы и текста задаёт шаблон (таблица~\ref{tab:page}). В документе их не меняют: не подключают пакеты, не задают шрифты, поля, колонтитулы и интервалы.

\begin{tabl}{|P{5cm}|L|}{Параметры страницы и текста}[tab:page]{\textbf{Параметр} & \textbf{Значение}}
Формат & A4, книжная ориентация \\ \hline
Поля & левое 25 мм, правое 15 мм, верхнее и нижнее 20 мм \\ \hline
Шрифт & Times New Roman, 12 пт. Если его нет на компьютере — Liberation Serif или CMU Serif \\ \hline
Межстрочный интервал & 1,15 \\ \hline
Абзацный отступ & 1,25 см, выравнивание по ширине \\ \hline
Верхний колонтитул & слева название документа, справа номер версии, под ними линия \\ \hline
Номер страницы & внизу по центру, на всех страницах кроме титульного листа \\ \hline
\end{tabl}

\section{Заголовки и нумерация}

Разделы нумеруются арабскими цифрами до трёх уровней: 1, 1.1, 1.1.1. Номера ставит шаблон, поэтому вручную номера в заголовки не пишут: иначе получится двойная нумерация. Точку в конце заголовка не ставят.

На раздел, таблицу, рисунок и требование ссылаются по номеру через \texttt{\textbackslash ref} и \texttt{\textbackslash reqref}. Номер, набранный вручную, после правок перестаёт совпадать с настоящим.

\section{Таблицы и рисунки}

\begin{itemize}
  \item Таблицы и рисунки нумеруются сквозной нумерацией по всему документу.
  \item Подпись таблицы стоит над ней, слева: «Таблица 1 – Название». Таблицу оформляют окружением \texttt{tabl}: при переходе на новую страницу шапка повторяется.
  \item Подпись рисунка стоит под ним, по центру: «Рисунок 1 – Название». Файлы рисунков лежат в \texttt{docs/img/}.
  \item На каждую таблицу и каждый рисунок в тексте есть ссылка, и она стоит до них.
\end{itemize}

\section{Язык и типографика}

\begin{itemize}
  \item Кавычки — «ёлочки». Диапазоны пишут коротким тире без пробелов: «стр. 75--76».
  \item Тире в тексте: \clarify{в «Требованиях к ПО» по замечанию проверяющего короткое «–», в остальных документах длинное «—»; нужно выбрать одно для всех документов}.
  \item Даты пишут в виде ДД.ММ.ГГГГ.
  \item Свои документы называют по-русски и в кавычках: «Системные требования», «План управления конфигурацией ПО». Пути к файлам, ссылки на README и на рабочие материалы команды в текст не пишут.
  \item На слайды и презентации курса, конспект лекций, гайд старших курсов и рабочие материалы команды не ссылаются: ни в списке ссылок, ни в тексте, ни в подписях, ни в определениях глоссария.
  \item Термины и сокращения берут из общего глоссария \texttt{docs/glossary.tex} и выводят командой \texttt{\textbackslash printterms}. Во всех документах термин определён одинаково.
\end{itemize}

\section{Запись требования}

Требование оформляют окружением \texttt{req}. Заголовок записи — идентификатор и версия: «SYS\_04\_02; \#1». Части идентификатора и сводную таблицу идентификаторов приводят в начале раздела с требованиями. Под текстом требования идут поля в постоянном порядке: «Источник», «Способ проверки», «Комментарий». Схема идентификаторов одна для всех документов, повтор идентификатора останавливает сборку.

Служебные символы «;~\#» и версия печатаются во всех документах: в заголовке требования («SRD\_01\_06; \#1») и в ссылке на родительское требование в поле «Источник» («SYS\_01\_05; \#1»). Требования объявляются по возрастанию идентификаторов. Проверку идентификации выполняют по разделу «Идентификация требований» чек-листа и командой \texttt{./build.sh -{}-check}.

\section{Выделение изменений}\label{sec:changes}

Когда документ сдают повторно, вместе с обычным PDF приносят экземпляр с выделенными изменениями. Изменения выделяются относительно принятой версии, а не относительно последнего сохранения. Поэтому в экземпляре видно всё, что изменилось с момента, когда преподаватель видел документ, даже если правки вносили несколько раз.

\begin{tabl}{|P{5cm}|L|}{Как выглядят изменения}[tab:diff]{\textbf{Что изменилось} & \textbf{Как показано}}
Добавленный текст & синим цветом с волнистым подчёркиванием \\ \hline
Удалённый текст & красным цветом, зачёркнут \\ \hline
Новая строка листа регистрации & синим цветом с волнистым подчёркиванием \\ \hline
Принятая версия & указана в рамке на титульном листе \\ \hline
\end{tabl}

Порядок работы:
\begin{enumerate}
  \item Принятые версии записаны в \texttt{docs/accepted.txt}: документ, коммит, номер версии, дата сдачи.
  \item Перед парой экземпляры собирают командой \texttt{./build.sh -{}-changes}. Они появляются в \texttt{pdf/changes/}.
  \item Когда преподаватель принял новую версию, в \texttt{docs/accepted.txt} заменяют коммит, номер версии и дату.
\end{enumerate}

Замена рисунка и изменение строения таблицы не выделяются. О них пишут в листе регистрации.

\section{Шаблоны документов}\label{sec:templates}

Документы одного типа строятся по одному шаблону (таблица~\ref{tab:templates}). Шаблон задаёт общий для всех документов раздел «Введение» и порядок остальных разделов. Новый документ начинают с копии шаблона: \texttt{docs/templates/<шаблон>.tex} копируют в \texttt{docs/0N\_<имя>.tex} и заменяют название, версию и дату.

\begin{tabl}{|P{3.8cm}|P{4.4cm}|L|}{Шаблоны документов}[tab:templates]{\textbf{Шаблон} & \textbf{Тип документа} & \textbf{Документы}}
\texttt{requirements.tex} & требования & «Системные требования», «Требования к ПО» \\ \hline
\texttt{plan.tex} & план & «План управления конфигурацией ПО», «План разработки ПО», в 8 семестре — «План верификации ПО» \\ \hline
\texttt{standard.tex} & стандарт & «Стандарт на разработку требований к ПО», в 8 семестре — стандарты на проектирование и кодирование \\ \hline
\end{tabl}

Раздел «Введение» одинаков во всех шаблонах: «Назначение», «Область применения», «Термины и сокращения», «Ссылки». Описание каждого процесса в планах и стандартах строится по одной схеме: цели, входные данные, условия входа, описание, условия выхода, выходные данные, роли.

\section{Проверка перед сдачей}\label{sec:check}

\begin{tabl}{|c|L|P{3.6cm}|}{Что проверяет автор}[tab:check]{\textbf{№} & \textbf{Проверка} & \textbf{Кто следит}}
1 & Документ собирается командой \texttt{./build.sh} без ошибок & шаблон \\ \hline
2 & Последняя строка листа регистрации совпадает с версией и датой на титульном листе & автор \\ \hline
3 & В заголовках нет номеров, набранных вручную & автор \\ \hline
4 & На каждую таблицу и каждый рисунок есть ссылка в тексте & автор \\ \hline
5 & Свои документы названы по-русски, в тексте нет путей к файлам и дат правок & автор \\ \hline
6 & Все красные пометки «[уточнить]» просмотрены & автор \\ \hline
7 & Для повторной сдачи собран экземпляр с выделенными изменениями & автор \\ \hline
\end{tabl}

\end{document}
